\subsection{Further uniform transverse sectors}

\begin{statement}{P4.44}{Proved}
In the remaining positive ruled case with generic transverse
type $A_3$ on every conductor component, the horizontal
section of the pole divisor is contracted by the resolution
map. The point-identification cycle occurs at a single
three-branch point of the critical curve.
\end{statement}
\begin{proof}
P4.42 excludes a cycle at two two-branch points. The
remaining graph has two rational normalized components
and one cycle, so its cycle occurs at one three-branch
point $p$. Suppose that the horizontal section $H$ is
not contracted. It maps finitely to one normalized
conductor component $C_H^\nu=\mathbb P^1$.
The generic tacnode has two conductor preimages.
The degree of $H\to\mathbb P^1$ is greater than one,
since $H$ has genus $g\ge1$, and hence is exactly two.
No other, vertical conductor component can map onto
$C_H$: the horizontal map already supplies both generic
preimages.

At a point where the entire critical curve is smooth,
P4R.14 gives either the product tacnode or the germ
$X(Y^2+XZ)$. Each has two distinct normalization points
over the critical-curve point. Each point lies on the
strict conductor, by P4R.6. Thus the two preimages must
both lie on $H$, and the degree-two map is unramified
there. All ramification is over singular points of $C$.

By (4.42a) and P4.33 the sum of the local defects is
$9-g\le8$. A nonplanar three-branch point has
$\delta\ge2$, and the branchwise estimate in P4.42
adds at least one per branch; its defect is at least
five. A planar three-branch point has $\delta\ge3$
and contributes at least $3\delta\ge9$, which is
already impossible. Every additional unibranch singular
point contributes at least two: in the planar case
the bound is $3\delta$, and in the nonplanar case it is
at least $\delta+1$. There is consequently at most one
additional unibranch singular point.

The normalization $C_H^\nu$ has at most two points
above $p$, since the other global conductor component
supplies at least one of its three branches. Including
the possible additional unibranch point leaves at most
three possible branch values for $H\to\mathbb P^1$.
Every ramification point of a degree-two map has index
two, and its branch value has only that ramified
preimage. Riemann--Hurwitz requires $2g+2\ge4$ distinct
branch values. This contradiction proves that $H$ is
contracted.
\end{proof}

\begin{statement}{P4.45}{Proved}
In an actual torus degeneration, a point of a planar
singular critical curve cannot have generic transverse
type $D_4$ along all its local branches.
\end{statement}
\begin{proof}
Write the full reduced critical-curve ideal as
$P=(u,h(x,y))$, where $h$ is reduced and singular.
Generic multiplicity three gives $f\in P^{(3)}$.
Powers of this complete-intersection prime-by-components
ideal are unmixed, so $P^{(3)}=P^3$. Thus there are
holomorphic coefficients such that
\[
 f=A u^3+B u^2h+Cuh^2+Dh^3.
\]
Introduce an independent variable $H$ and the actual
holomorphic function
\[
 F(u,x,y,H)=A(u,x,y)u^3+B(u,x,y)u^2H
               +C(u,x,y)uH^2+D(u,x,y)H^3.
\]
Then $f=F(u,x,y,h)$. Arbitrarily small constant changes
of the four coefficients preserve the $D_4$ transverse
germs on a fixed boundary and make the binary cubic
at the origin have distinct roots and a simple critical
point as a polynomial in $u/H$. Its critical value is
nonzero: a zero critical value would be a repeated root.
We also choose its leading coefficient nonzero.

After substituting $u=Hv$, the quotient $F_u/H^2$
is holomorphic. The implicit-function theorem gives
a holomorphic critical root $v=v(x,y,H)$ and
\[
          F(Hv(x,y,H),x,y,H)=H^3\beta(x,y,H),
                       \qquad \beta(0)\ne0.
\]
Choose the local cube root of the unit $\beta$.
The map $H\mapsto H\beta(x,y,H)^{1/3}$ has nonzero
$H$ derivative at the origin and has a holomorphic
inverse with parameters $x,y$. Put
\[
             \widehat h(x,y)=h(x,y)
                           \beta(x,y,h(x,y))^{1/3}.
\]
It defines the same singular reduced plane curve as $h$.
Choose a small morsification $\widehat h_s$ with a
Morse critical point of nonzero critical value. This
exists by the positive local degree of the gradient
of the singular reduced plane function; a small constant
shift avoids zero critical values. Use the displayed
inverse to solve exactly
\[
         H_s\beta(x,y,H_s)^{1/3}=\widehat h_s,
             \qquad f_s=F(u,x,y,H_s(x,y)).
\]
On its $u$-critical branch the critical-value function
is exactly $\widehat h_s^3$. At the selected plane
Morse point, $H_s\ne0$ and $F_{uu}\ne0$. The Schur
complement of the full Hessian is
$3\widehat h_s^2\operatorname{Hess}(\widehat h_s)$.
Consequently $f_s$ has a nondegenerate three-variable
critical point with nonzero critical value.

All deformations above can be chosen successively small
on the fixed boundary. There the moving curve $H_s=0$
has the same ordinary-triple transverse germ: its binary
cubic has three distinct linear factors. The boundary
transport used in P4.28 therefore identifies the local
nearby fibres through nested representatives. The Morse
vanishing sphere embeds in a nearby global torus fibre
and has self-intersection $-2$, contradicting its
asphericity. This proves the assertion.
\end{proof}

\begin{statement}{P4.46}{Proved}
In the requested full-fibration classes, a support with
positive-genus ruled resolution cannot have positive
Euler number and generic transverse type $D_4$ along
every component of its critical curve.
\end{statement}
\begin{proof}
If the horizontal section is contracted, P4.39 already
gives a contradiction. Suppose it is not contracted.
The remaining critical curve has two rational normalized
components and one graph cycle. For the actual
conductor scheme, ordinary-triple transverse type gives
the ideal $P^{(2)}$, as in P4.31. Let $Q=P/P^{(2)}$,
and let $\tau_p$ be the full colength of its differentiated
image in the direct image of the normalized rank-two
kernel $K$. This includes both the branchwise lattice
colength and the further descent colength. Riemann--Roch
on the two normalized projective lines gives
\[
 \chi(\mathcal O_{C_{\rm sch}})
       =10-\sum_p(\delta_p+\tau_p-r_p).
\]
P4.33 identifies the left side with $g-1$, so the total
defect is $11-g\le10$.

We prove that a three-branch point contributes at least
eleven. A planar point contributes $5\delta\ge15$ by
the exact different calculation in P4.34. Consider a
nonplanar point. If its three branches are smooth, the
curve has multiplicity three and embedding dimension
three. A generic linear parameter has an Artinian
section of length three and embedding dimension two,
with ideal $(v,w)^2$. The defining curve ideal therefore
has exactly three minimal generators. To justify the
lifting, the parameter is a nonzerodivisor on the
Cohen--Macaulay curve ring, so the Tor term vanishes in
the ideal sequence modulo that parameter; Nakayama
then lifts precisely the three generators.
Each branchwise conormal lattice has rank two and
colength at least two, since every defining gradient
has positive order. Their sum has colength at least
six. Modulo the three branch maximal ideals its
normalization has six generator values. The image of
$Q$, generated by three elements over the local curve
ring, supplies a vector space of dimension at most
three. The additional descent colength is therefore
at least three. Since $\delta\ge2$ and $R=0$, the
full defect is at least $2+6+3=11$.

If one of the three branches is singular, let $m\ge4$
be total multiplicity. The branchwise gradient bound
gives $\tau\ge2m$, while $\deg R=m-3$.
The general inequality $\delta\ge m-1$ leaves only
$m=4,\delta=3$ as a possible value below the required
bound. That value has excess $\delta-(3-1)=1$.
The exact classification in P4.40 shows that its only
singular-branch model is a cusp and two independent
lines, of embedding dimension four. It does not embed
in the ambient smooth threefold. Thus $\delta\ge4$
in the present case, and the full defect is at least
$\delta+2m-(m-3)=\delta+m+3\ge11$.
A cycle at one three-branch point is impossible.

The other possible cycle has two two-branch points.
Each contributes at least five by P4.34; equality is
possible only for a planar ordinary node. A nonplanar
two-branch point contributes at least seven, and a
nonordinary planar point at least ten. The total bound
of ten therefore forces two ordinary planar nodes
and $g=1$. Each such node is forbidden by P4.45.
This proves the assertion.
\end{proof}

The local Hessian identities, rank-one lattice identity,
and exact integer defect bounds are checked in
\texttt{code/part4/uniform\_sectors.py}; the saved output
reports PASS.
